\documentclass[10pt,a4paper]{amsart}
\usepackage[margin=19mm]{geometry}
\usepackage{amsmath,amssymb,mathtools}
\usepackage[scr=rsfs,cal=euler]{mathalfa}
\usepackage{xfrac}
\usepackage[unicode,hidelinks,bookmarks=false]{hyperref}
\newcommand{\ie}{i.e.\ }
\newcommand{\cf}{cf.\ }
\newcommand{\resp}{respectively\ }
\newcommand{\Subsection}{\subsection}
\providecommand{\nobreakdash}{\nobreak\hbox{-}}
\setlength{\emergencystretch}{2em}
\multlinegap0pt
\usepackage{amssymb}
\usepackage{mathtools}
\newtheorem{lemma}{Lemma}
\newtheorem{theorem}{Theorem}
\title{A viscosity-Halpern hybrid scheme for countable families of equilibrium and variational inequality problems}
\author{Markjoe O. Uba}
\date{Source: arXiv:2606.03088v1 (CC BY 4.0)}
\begin{document}
\maketitle

\noindent\emph{Source and licence.} A viscosity-Halpern hybrid scheme for countable families of equilibrium and variational inequality problems. Version arXiv:2606.03088v1 (CC BY 4.0), distributed by arXiv under the Creative Commons Attribution 4.0 International licence, http://creativecommons.org/licenses/by/4.0/, as recorded in the arXiv licence metadata for this exact version and shown on the abstract page https://arxiv.org/abs/2606.03088v1. Full licence notice: \texttt{licenses/arxiv-2606.03088-CC-BY-4.0.txt}.\par\smallskip

\noindent\textit{Editorial scope.} Counted unit: the article's theorem mainthm (source0.tex lines 300--310). The article's complete original proof of it is delivered with it (source0.tex lines 312--575, one \texttt{proof} environment of 264 lines). No mathematics is written, repaired or re-derived: the statement and the proof are the article's own lines, in the article's own notation, and no retained line is substituted or re-flowed.  Retained uncounted as the source-local context the proof needs: lines 197--202; lines 204--209; lines 211--222; lines 224--233; lines 235--244; lines 264--268; lines 274--285.  Nothing was omitted from any retained span, so the package carries no omission record.  No retained line contains a citation, so the wrapper carries no bibliography and no bibliography entry.  No retained line contains a figure environment, an image-inclusion command or drawing code, so no image is delivered and the package declares no asset dependency.  Editorial formatting only, in addition: an AMS \texttt{amsart} preamble that restates the article's own declarations and macros used by the retained lines, namely the environments \texttt{lemma}, \texttt{theorem} on separate counters, together with the standard packages those lines need.  The retained environments print in this document as Lemma 1, Theorem 1 in order of appearance, whereas the article prints the same items with the numbering of its own full document.  The complete native source of the article as distributed in the arXiv e-print for this exact version is bundled unchanged as \texttt{sources/201956/source0.tex}.
\begin{lemma}\label{glemma}
Let $E$ be a uniformly convex Banach space, $r>0$ and $B_r(0)$ be the closed ball of $E$. For any given points $x_1,x_2,\ldots,x_N\in B_r(0)$ and any positive numbers $\lambda_1,\lambda_2,\ldots,\lambda_N$ with $\sum_{i=1}^{N}\lambda_i=1$, there exists a continuous, strictly increasing and convex function $g:[0,2r)\to [0,\infty)$ with $g(0)=0$ such that, for any $i,j\in\{1,2,\ldots,N\}$ with $i<j$,
\[
\left\|\sum_{n=1}^{N}\lambda_nx_n\right\|^2\leq \sum_{n=1}^{N}\lambda_n\|x_n\|^2-\lambda_i\lambda_jg(\|x_i-x_j\|).
\]
\end{lemma}

\begin{lemma}\label{phitozero}
Let $E$ be a real smooth and uniformly convex Banach space, and let $\{x_n\}$ and $\{y_n\}$ be two sequences in $E$. If either $\{x_n\}$ or $\{y_n\}$ is bounded and $\phi(x_n,y_n)\to 0$ as $n\to\infty$, then
\[
\|x_n-y_n\|\to 0.
\]
\end{lemma}

\begin{lemma}[Generalized projection inequality]\label{projineq}
Let $E$ be a smooth, strictly convex and reflexive Banach space, and let $C$ be a nonempty closed and convex subset of $E$. Then $z=\Pi_Cx$ if and only if
\[
\langle y-z,Jx-Jz\rangle\leq 0,
\qquad \forall y\in C.
\]
Moreover,
\[
\phi(y,z)+\phi(z,x)\leq \phi(y,x),
\qquad \forall y\in C.
\]
\end{lemma}

\begin{lemma}\label{EPrezolvent}
Let $C$ be a nonempty closed subset of a smooth, strictly convex and reflexive Banach space $E$ such that $JC$ is closed and convex. Let $f$ be a bifunction from $JC\times JC$ to $\mathbb{R}$ satisfying $(A1)-(A4)$. For $r>0$ and $x\in E$, define $T_r^f:E\to C$ by
\[
T_r^f x=\left\{z\in C:f(Jz,Jy)+\frac{1}{r}\langle y-z,Jz-Jx\rangle\geq 0,\ \forall y\in C\right\}.
\]
Then $T_r^f$ is single valued, $F(T_r^f)=EP(f)$, $JEP(f)$ is closed and convex, and
\[
\phi(p,T_r^f x)+\phi(T_r^f x,x)\leq \phi(p,x),\qquad \forall p\in EP(f).
\]
\end{lemma}

\begin{lemma}\label{VIrezolvent}
Let $C$ be a nonempty closed subset of a smooth, strictly convex and reflexive Banach space $E$. Let $A:C\to E^{*}$ be a continuous monotone mapping. For $r>0$ and $x\in E$, define $F_r^A:E\to C$ by
\[
F_r^A x=\left\{z\in C:\langle y-z,Az\rangle+\frac{1}{r}\langle y-z,Jz-Jx\rangle\geq 0,\ \forall y\in C\right\}.
\]
Then $F_r^A$ is single valued, $F(F_r^A)=VI(C,A)$, $JVI(C,A)$ is closed and convex, and
\[
\phi(p,F_r^A x)+\phi(F_r^A x,x)\leq \phi(p,x),\qquad \forall p\in VI(C,A).
\]
\end{lemma}

\begin{equation}\label{lambdacond}
\lambda_n+\sigma_n<1,
\qquad \lambda_n\to 0,
\qquad \sigma_n\to 0.
\end{equation}

\begin{equation}\label{algorithm}
\left\{
\begin{array}{ll}
s_n=J^{-1}\big(\lambda_nJh(x_n)+\sigma_nJ\bar u+(1-\lambda_n-\sigma_n)Jx_n\big),\\[2mm]
z_n=\left\{z\in C:f_n(Jz,Jy)+\frac{1}{r_n}\langle y-z,Jz-Js_n\rangle\geq 0,\ \forall y\in C\right\},\\[2mm]
u_n=\left\{z\in C:\langle y-z,A_nz\rangle+\frac{1}{r_n}\langle y-z,Jz-Js_n\rangle\geq 0,\ \forall y\in C\right\},\\[2mm]
y_n=J^{-1}\big(\alpha_1Js_n+\alpha_2Jz_n+\alpha_3T_nu_n\big),\\[2mm]
C_{n+1}=\{v\in C_n:\phi(v,y_n)\leq \phi(v,s_n)\},\\[2mm]
x_{n+1}=\Pi_{C_{n+1}}x_1,
\end{array}
\right.
\end{equation}

\begin{theorem}\label{mainthm}
Let $E$ be a uniformly smooth and uniformly convex real Banach space with dual space $E^{*}$, and let $C$ be a nonempty closed and convex subset of $E$ such that $JC$ is closed and convex. Let $\{f_i\}_{i=1}^{\infty}$ be a countable family of bifunctions from $JC\times JC$ to $\mathbb{R}$ satisfying $(A1)-(A4)$, let $\{A_j\}_{j=1}^{\infty}$ be a countable family of continuous monotone mappings from $C$ into $E^{*}$, and let $\{T_n\}_{n=1}^{\infty}$ be a countable family of generalized $J_{*}$-nonexpansive maps. Let $\Gamma$ be a family of $J_{*}$-closed and generalized $J_{*}$-nonexpansive maps from $C$ into $E^{*}$ such that
\[
\bigcap_{n=1}^{\infty}F_J(T_n)=F_J(\Gamma)\neq\emptyset.
\]
Assume that $\{T_n\}$ satisfies the NST-condition with $\Gamma$, and suppose that the common solution set
\[
B:=F_J(\Gamma)\cap\left[\bigcap_{i=1}^{\infty}EP(f_i)\right]\cap\left[\bigcap_{j=1}^{\infty}VI(C,A_j)\right]
\]
is nonempty, closed and convex. Then the sequence $\{x_n\}$ generated by \eqref{algorithm} is well defined and converges strongly to $\Pi_Bx_1$, the generalized projection of $x_1$ onto $B$.
\end{theorem}

\begin{proof}
We divide the proof into seven steps.

\noindent {\bf Step 1.} We show that the construction is well defined and that $B\subset C_n$ for all $n\geq 1$.

Clearly, $B\subset C_1=C$. Suppose that $B\subset C_n$ for some $n\geq 1$. Let $p\in B$. Then $p\in EP(f_n)$ and $p\in VI(C,A_n)$. By Lemmas~\ref{EPrezolvent} and \ref{VIrezolvent}, and by the definitions of $z_n$ and $u_n$, we have
\begin{equation}\label{resolventineq}
\phi(p,z_n)\leq \phi(p,s_n),\qquad \phi(p,u_n)\leq \phi(p,s_n).
\end{equation}
Since $p\in F_J(T_n)$ and $T_n$ is generalized $J_{*}$-nonexpansive,
\begin{equation}\label{Tineq}
\phi(p,(J_{*}\circ T_n)u_n)\leq \phi(p,u_n)\leq \phi(p,s_n).
\end{equation}
Using the definition of $y_n$ and the convexity inequality for the square of the norm in $E^{*}$, we obtain
\begin{align*}
\phi(p,y_n)
&=\phi\big(p,J^{-1}(\alpha_1Js_n+\alpha_2Jz_n+\alpha_3T_nu_n)\big)\\
&\leq \alpha_1\phi(p,s_n)+\alpha_2\phi(p,z_n)+\alpha_3\phi(p,(J_{*}\circ T_n)u_n)\\
&\leq \phi(p,s_n).
\end{align*}
Thus $p\in C_{n+1}$, and hence $B\subset C_{n+1}$.

Moreover, $C_{n+1}$ is closed and convex because the inequality
\[
\phi(v,y_n)\leq \phi(v,s_n)
\]
is equivalent to
\[
2\langle v,Js_n-Jy_n\rangle\leq \|s_n\|^2-\|y_n\|^2,
\]
which defines a closed half-space intersected with $C_n$. Since $B\subset C_{n+1}$, the set $C_{n+1}$ is nonempty. By induction, each $C_n$ is nonempty, closed and convex. Therefore $\Pi_{C_n}x_1$ exists for each $n$, and the algorithm is well defined.

\noindent {\bf Step 2.} We show that $\{x_n\}$ converges strongly to some point $x^{*}\in C$.

Since $x_n=\Pi_{C_n}x_1$ and $B\subset C_n$, we have
\[
\phi(x_n,x_1)\leq \phi(p,x_1),\qquad \forall p\in B.
\]
Thus $\{\phi(x_n,x_1)\}$ is bounded, and consequently $\{x_n\}$ is bounded. Also, since $x_{n+1}\in C_{n+1}\subset C_n$ and $x_n=\Pi_{C_n}x_1$, we have
\[
\phi(x_n,x_1)\leq \phi(x_{n+1},x_1),\qquad n\geq 1.
\]
Therefore $\lim_{n\to\infty}\phi(x_n,x_1)$ exists.

Let $m>n$. Since $x_m\in C_m\subset C_n$, Lemma~\ref{projineq} gives
\begin{align}\label{Cauchyineq}
\phi(x_m,x_n)&\leq \phi(x_m,x_1)-\phi(x_n,x_1)\to 0
\end{align}
as $m,n\to\infty$. By Lemma~\ref{phitozero}, $\|x_m-x_n\|\to 0$ as $m,n\to\infty$. Hence $\{x_n\}$ is Cauchy. Since $C$ is closed, there exists $x^{*}\in C$ such that
\begin{equation}\label{xnconv}
x_n\to x^{*}.
\end{equation}

\noindent {\bf Step 3.} We show that $s_n\to x^{*}$ and $y_n\to x^{*}$.

Since $\{x_n\}$ is bounded and $h$ is a contraction, the sequence $\{h(x_n)\}$ is bounded. Also $\bar u$ is fixed. By \eqref{lambdacond},
\[
\lambda_nJh(x_n)+\sigma_nJ\bar u+(1-\lambda_n-\sigma_n)Jx_n-Jx_n\to 0.
\]
Since $J^{-1}=J_{*}$ is uniformly continuous on bounded subsets of $E^{*}$, we obtain
\begin{equation}\label{snconv}
\|s_n-x_n\|\to 0.
\end{equation}
Combining \eqref{xnconv} and \eqref{snconv}, we get
\begin{equation}\label{snxstar}
s_n\to x^{*}.
\end{equation}

Since $x_{n+1}\in C_{n+1}$, we have
\[
\phi(x_{n+1},y_n)\leq \phi(x_{n+1},s_n).
\]
Using \eqref{xnconv} and \eqref{snxstar}, we obtain $\phi(x_{n+1},s_n)\to 0$. Hence $\phi(x_{n+1},y_n)\to 0$. By Lemma~\ref{phitozero},
\[
\|x_{n+1}-y_n\|\to 0.
\]
Since $x_{n+1}\to x^{*}$, it follows that
\begin{equation}\label{ynconv}
y_n\to x^{*}.
\end{equation}

\noindent {\bf Step 4.} We show that $u_n\to x^{*}$, $z_n\to x^{*}$ and $\|Ju_n-T_nu_n\|\to 0$.

Let $p\in B$. The sequences $\{Js_n\}$, $\{Jz_n\}$ and $\{T_nu_n\}$ are bounded in $E^{*}$. Applying Lemma~\ref{glemma} in $E^{*}$ to the bounded ball containing these points, we obtain a continuous, strictly increasing and convex function $g$ with $g(0)=0$ such that
\begin{align}\label{gineq}
\phi(p,y_n)
&\leq \alpha_1\phi(p,s_n)+\alpha_2\phi(p,z_n)+\alpha_3\phi(p,(J_{*}\circ T_n)u_n)\nonumber\\
&\quad -\alpha_1\alpha_3g(\|Js_n-T_nu_n\|)\nonumber\\
&\leq \phi(p,s_n)-\alpha_1\alpha_3g(\|Js_n-T_nu_n\|).
\end{align}
Since $s_n\to x^{*}$ and $y_n\to x^{*}$, we have
\[
\phi(p,s_n)-\phi(p,y_n)\to 0.
\]
It follows from \eqref{gineq} that
\[
g(\|Js_n-T_nu_n\|)\to 0.
\]
Since $g$ is strictly increasing and $g(0)=0$, we get
\begin{equation}\label{residual1}
\|Js_n-T_nu_n\|\to 0.
\end{equation}

Furthermore, from \eqref{gineq} and the fact that each term in the convex combination is not larger than $\phi(p,s_n)$, we obtain
\begin{equation}\label{conv_phi_terms}
\phi(p,z_n)\to \phi(p,x^{*}),\qquad \phi(p,u_n)\to \phi(p,x^{*}).
\end{equation}
Using Lemmas~\ref{EPrezolvent} and \ref{VIrezolvent}, we have
\[
\phi(z_n,s_n)\leq \phi(p,s_n)-\phi(p,z_n),
\]
and
\[
\phi(u_n,s_n)\leq \phi(p,s_n)-\phi(p,u_n).
\]
By \eqref{snxstar} and \eqref{conv_phi_terms}, the right hand sides tend to zero. Thus
\[
\phi(z_n,s_n)\to 0,
\qquad
\phi(u_n,s_n)\to 0.
\]
By Lemma~\ref{phitozero},
\begin{equation}\label{unznsn}
\|z_n-s_n\|\to 0,
\qquad
\|u_n-s_n\|\to 0.
\end{equation}
Combining \eqref{snxstar} and \eqref{unznsn}, we obtain
\begin{equation}\label{unzconv}
z_n\to x^{*},
\qquad
u_n\to x^{*}.
\end{equation}
Since $J$ is uniformly continuous on bounded subsets of $E$, \eqref{unznsn} implies
\[
\|Ju_n-Js_n\|\to 0.
\]
Together with \eqref{residual1}, we get
\begin{equation}\label{residual2}
\|Ju_n-T_nu_n\|\leq \|Ju_n-Js_n\|+\|Js_n-T_nu_n\|\to 0.
\end{equation}

\noindent {\bf Step 5.} We prove that $x^{*}\in F_J(\Gamma)$.

From \eqref{unzconv}, $u_n\to x^{*}$. From \eqref{residual2},
\[
\|Ju_n-T_nu_n\|\to 0.
\]
Since $\{T_n\}$ satisfies the NST-condition with $\Gamma$, we obtain
\begin{equation}\label{NSTuse}
\|Ju_n-Tu_n\|\to 0,
\qquad \forall T\in \Gamma.
\end{equation}
Let $T\in\Gamma$. Since $u_n\to x^{*}$ and $T$ is $J_{*}$-closed, \eqref{NSTuse} implies that $Tx^{*}=Jx^{*}$. Hence $x^{*}\in F_J(T)$ for every $T\in\Gamma$. Therefore
\begin{equation}\label{fixedmembership}
x^{*}\in F_J(\Gamma).
\end{equation}

\noindent {\bf Step 6.} We prove that
\[
x^{*}\in \left[\bigcap_{i=1}^{\infty}EP(f_i)\right]\cap\left[\bigcap_{j=1}^{\infty}VI(C,A_j)\right].
\]

We first show that $x^{*}\in \bigcap_{j=1}^{\infty}VI(C,A_j)$. From the definition of $u_n$,
\begin{equation}\label{VIbasic}
\langle y-u_n,A_nu_n\rangle+\frac{1}{r_n}\langle y-u_n,Ju_n-Js_n\rangle\geq 0,
\qquad \forall y\in C.
\end{equation}
Since $r_n\geq a>0$ and $\|Ju_n-Js_n\|\to 0$, we have
\begin{equation}\label{VIfraction}
\frac{\|Ju_n-Js_n\|}{r_n}\to 0.
\end{equation}
Fix $j\in\mathbb{N}$. Since $\nu^{-1}(j)$ is infinite, choose a subsequence $\{n_m\}$ such that $A_{n_m}=A_j$ for all $m\geq 1$. From \eqref{VIbasic},
\[
\langle y-u_{n_m},A_ju_{n_m}\rangle+\frac{1}{r_{n_m}}\langle y-u_{n_m},Ju_{n_m}-Js_{n_m}\rangle\geq 0,
\qquad \forall y\in C.
\]
For $t\in(0,1]$ and $y\in C$, set $v_t=ty+(1-t)x^{*}$. Since $C$ is convex, $v_t\in C$. Using monotonicity of $A_j$, we obtain
\[
\langle v_t-u_{n_m},A_jv_t\rangle\geq -\frac{1}{r_{n_m}}\langle v_t-u_{n_m},Ju_{n_m}-Js_{n_m}\rangle.
\]
Letting $m\to\infty$ and using \eqref{unzconv} and \eqref{VIfraction}, we get
\[
\langle v_t-x^{*},A_jv_t\rangle\geq 0.
\]
Since $v_t-x^{*}=t(y-x^{*})$, it follows that
\[
\langle y-x^{*},A_jv_t\rangle\geq 0.
\]
Letting $t\downarrow 0$ and using the continuity of $A_j$, we obtain
\[
\langle y-x^{*},A_jx^{*}\rangle\geq 0,
\qquad \forall y\in C.
\]
Thus $x^{*}\in VI(C,A_j)$. Since $j\in\mathbb{N}$ was arbitrary,
\begin{equation}\label{VIall}
x^{*}\in \bigcap_{j=1}^{\infty}VI(C,A_j).
\end{equation}

We now show that $x^{*}\in \bigcap_{i=1}^{\infty}EP(f_i)$. From the definition of $z_n$,
\begin{equation}\label{EPbasic}
f_n(Jz_n,Jy)+\frac{1}{r_n}\langle y-z_n,Jz_n-Js_n\rangle\geq 0,
\qquad \forall y\in C.
\end{equation}
Since $r_n\geq a>0$ and $\|Jz_n-Js_n\|\to 0$, we have
\begin{equation}\label{EPfraction}
\frac{\|Jz_n-Js_n\|}{r_n}\to 0.
\end{equation}
Fix $i\in\mathbb{N}$. Since $\mu^{-1}(i)$ is infinite, choose a subsequence $\{m_q\}$ such that $f_{m_q}=f_i$ for all $q\geq 1$. From \eqref{EPbasic} and the monotonicity condition $(A2)$,
\[
\frac{1}{r_{m_q}}\langle y-z_{m_q},Jz_{m_q}-Js_{m_q}\rangle\geq -f_i(Jz_{m_q},Jy)\geq f_i(Jy,Jz_{m_q}).
\]
Letting $q\to\infty$ and using \eqref{unzconv}, \eqref{EPfraction} and the lower semicontinuity of $f_i(Jy,\cdot)$, we obtain
\[
f_i(Jy,Jx^{*})\leq 0,
\qquad \forall y\in C.
\]
For $t\in(0,1]$ and $y\in C$, put
\[
y_t^{*}=tJy+(1-t)Jx^{*}.
\]
Since $JC$ is convex, $y_t^{*}\in JC$. Hence $f_i(y_t^{*},Jx^{*})\leq 0$. By $(A1)$ and the convexity of the second argument,
\[
0=f_i(y_t^{*},y_t^{*})\leq t f_i(y_t^{*},Jy)+(1-t)f_i(y_t^{*},Jx^{*})\leq t f_i(y_t^{*},Jy).
\]
Thus $f_i(y_t^{*},Jy)\geq 0$. Letting $t\downarrow 0$ and using $(A3)$, we obtain
\[
f_i(Jx^{*},Jy)\geq 0,
\qquad \forall y\in C.
\]
Therefore $x^{*}\in EP(f_i)$. Since $i\in\mathbb{N}$ was arbitrary,
\begin{equation}\label{EPall}
x^{*}\in \bigcap_{i=1}^{\infty}EP(f_i).
\end{equation}
Combining \eqref{fixedmembership}, \eqref{VIall} and \eqref{EPall}, we have
\begin{equation}\label{xinstarB}
x^{*}\in B.
\end{equation}

\noindent {\bf Step 7.} We show that $x^{*}=\Pi_Bx_1$.

Since $B$ is nonempty, closed and convex, $\Pi_Bx_1$ exists. By the defining property of the generalized projection,
\begin{equation}\label{PBineq1}
\phi(\Pi_Bx_1,x_1)\leq \phi(x^{*},x_1).
\end{equation}
Since $B\subset C_n$ for every $n\geq 1$ and $x_n=\Pi_{C_n}x_1$, we also have
\[
\phi(x_n,x_1)\leq \phi(\Pi_Bx_1,x_1),
\qquad n\geq 1.
\]
Letting $n\to\infty$, we obtain
\begin{equation}\label{PBineq2}
\phi(x^{*},x_1)\leq \phi(\Pi_Bx_1,x_1).
\end{equation}
From \eqref{PBineq1} and \eqref{PBineq2},
\[
\phi(x^{*},x_1)=\phi(\Pi_Bx_1,x_1).
\]
Since $x^{*}\in B$ and the generalized projection is unique, we conclude that
\[
x^{*}=\Pi_Bx_1.
\]
This completes the proof.
\end{proof}

\end{document}
